\chapter*{ЗАКЛЮЧЕНИЕ}
\addcontentsline{toc}{chapter}{ЗАКЛЮЧЕНИЕ}

В работе рассмотрена проблема токенизации в геномных языковых моделях и предложен метод, который адаптирует границы токенов к локальной структуре ДНК. В отличие от статических схем (посимвольной, $k$-мерной и BPE), предложенная архитектура обучает механизм сегментации совместно с языковой моделью и тем самым связывает качество токенизации с целевой задачей предсказания следующего нуклеотида.

В первой главе показано, что современные DNA~LM достигли существенного прогресса, однако вопрос токенизации остаётся открытым: статические правила разбиения плохо согласуются с функциональной неоднородностью генома. При этом недавние работы по обучаемой токенизации подтверждают интерес к направлению, но обычно закрывают только часть задачи - например, дают адаптивность без биологических приоров или используют биоприоры без совместного end-to-end обучения.

Во второй главе формализована архитектура адаптивного токенизатора на основе рекурсивного гейтированного слияния, дифференцируемого динамического программирования и каузального декодирования обратно в нуклеотидное разрешение. Дополнительно описана интеграция биологических приоров (phastCons и JASPAR), составная функция потерь и трёхфазная стратегия обучения, которая уменьшает разрыв между мягкими решениями на обучении и жёсткими на инференсе.

Итоговый вклад работы состоит в связной постановке задачи и в конкретной, воспроизводимой схеме её решения: от математической формулировки выбора границ до плана сравнительной оценки на DART-Eval и смежных геномных бенчмарках. Практически это означает, что предложенный модуль можно встроить в существующие пайплайны предобучения DNA~LM без пересборки всей архитектуры с нуля.

В качестве продолжения логично выполнить полный экспериментальный цикл: предобучение с абляциями, проверку устойчивости на нескольких референсах генома и анализ того, где именно модель предпочитает короткие или длинные спаны. По этим результатам уже можно сделать окончательный вывод о том, насколько адаптивная токенизация улучшает downstream-задачи регуляторной геномики в сравнении с сильными статическими и гибридными базовыми решениями.
